\section{Notation and Problem Statement}
\label{sec:problem}

This section fixes the notation used in the rest of the paper and states the problem that the
method solves. Every symbol is defined here before it is used anywhere else.

\subsection{Notation}
\label{sec:notation}
Images are real valued arrays of height $H$ and width $W$ with values in $[0,1]$. We write $\mathbf{y}$
for the noisy B scan that comes off the scanner and $\mathbf{x}$ for the clean reference obtained by
averaging many repeated scans of the same location. A pixel index is written $p$, so $y_p$ is one
pixel of $\mathbf{y}$.

Two products appear. The symbol $\odot$ is the Hadamard product, which multiplies two arrays of the
same shape entry by entry, so $(\mathbf{a} \odot \mathbf{b})_p = a_p b_p$. Ordinary multiplication of
two numbers is written by juxtaposition with no symbol between the factors. Table~\ref{tab:notation}
lists every remaining symbol in the order in which the paper introduces it.

\begin{table}[!tb]
\centering
\caption{Symbols used in the paper, in order of first use.}
\label{tab:notation}
\setlength{\tabcolsep}{4pt}
\footnotesize
\begin{tabular}{lp{0.62\columnwidth}}
\toprule
Symbol & Meaning \\
\midrule
$\mathbf{y}$ & Noisy input B scan \\
$\mathbf{x}$ & Clean reference B scan \\
$f_\theta$ & Frozen denoising backbone with parameters $\theta$ \\
$\mathbf{b}$ & Backbone output, $\mathbf{b} = f_\theta(\mathbf{y})$ \\
$\phi$ & Trainable parameters of the wrapper \\
$\mathbf{z}$ & Backbone feature map read by the confidence head \\
$\mathbf{c}$ & Confidence map with values in $[0,1]$ \\
$\mathbf{u}$ & Uncertainty map, $\mathbf{u} = \mathbf{1} - \mathbf{c}$ \\
$s_i$ & Score of clinical property $i$, a number in $[0,1]$ \\
$\mathbf{m}_i$ & Failure map of clinical property $i$, values in $[0,1]$ \\
$\tau_i$ & Decision threshold for property $i$ \\
$f_i$ & Soft failure signal derived from $s_i$ and $\tau_i$ \\
$\boldsymbol{\Delta}$ & Raw multiplicative gain proposed by the gain network \\
$\boldsymbol{\lambda}$ & Per pixel strength map with values in $[0,1]$ \\
$\mathbf{a}$ & Allocation map produced by the rule layer \\
$\boldsymbol{\alpha}$ & Gain after allocation, strength and clamping \\
$\mathbf{g}$ & Intensity gate that suppresses change in the background \\
$\boldsymbol{\delta}_{\mathrm{e}}$ & Additive edge correction \\
$\mathbf{q}$ & Candidate image before the safety decision \\
$n$ & Number of satisfied safety constraints, an integer in $\{0,1,2,3\}$ \\
$w$ & Blend weight chosen from $n$ \\
$\hat{\mathbf{x}}$ & Final output of the wrapper \\
$\mathcal{T},\mathcal{B}$ & Sets of tissue pixels and background pixels \\
$\mu_\mathcal{T},\mu_\mathcal{B}$ & Mean value over tissue and over background \\
$\sigma_\mathcal{B}$ & Standard deviation over background \\
\bottomrule
\end{tabular}
\end{table}

\subsection{The optimisation problem}
The denoiser is given and is not changed. The task is to produce an image that beats
$\mathbf{b} = f_\theta(\mathbf{y})$ on named clinical measures, stays close to it everywhere, and can
explain every change it makes.
\label{sec:optimisation}
We now state the problem in the form that the training procedure actually solves. This makes precise
what is being optimised, over what, and subject to what.

\textbf{Decision variables.} The only free quantities are the wrapper parameters $\phi$. These are the
confidence head, the gain network, the edge network, the strength predictor and the rule layer. The
backbone parameters $\theta$ are held fixed throughout.

\textbf{Objective.} Let $\mathcal{D}$ be a set of paired scans. Let $\hat{\mathbf{x}}(\phi)$ be the
output of the wrapper. We minimise the expected value of a weighted sum of five terms,
\begin{equation}
\min_{\phi}\;\; \mathbb{E}_{(\mathbf{y},\mathbf{x}) \sim \mathcal{D}}
\Big[\, \mathcal{L}_{\mathrm{fid}} + w_c \mathcal{L}_{\mathrm{clin}}
      + w_p \mathcal{L}_{\mathrm{pred}} + w_o \mathcal{L}_{\mathrm{coop}}
      + w_s \mathcal{L}_{\mathrm{edge}} \,\Big],
\label{eq:objective}
\end{equation}
where $\mathcal{L}_{\mathrm{fid}}$ keeps the output near the clean reference,
$\mathcal{L}_{\mathrm{clin}}$ rewards improvement of the clinical measures over the backbone,
$\mathcal{L}_{\mathrm{pred}}$ calibrates the property scores against their thresholds,
$\mathcal{L}_{\mathrm{coop}}$ ties the amount of correction to the estimated uncertainty, and
$\mathcal{L}_{\mathrm{edge}}$ supervises the additive edge branch. Section~\ref{sec:training} gives
each term in full and lists the weights.

\textbf{Constraints.} Three requirements must hold for the output on every image. They are stated
formally in Section~\ref{sec:safety} and summarised here. The aggregate clinical quality must not
fall by more than a fixed tolerance. No single clinical property may fall by more than a fixed
tolerance while another rises. The relative change from the backbone output must stay below a fixed
bound.

\textbf{How the problem is solved.} The objective is differentiable in $\phi$, including the rule
layer, because every logical connective is replaced by a smooth piecewise linear surrogate. We
therefore solve the problem by mini batch stochastic gradient descent with the AdamW
optimiser~\cite{loshchilov2019decoupled}. The three constraints are handled in two ways. During
training they enter the objective as penalties, so the optimiser is pushed toward feasible
solutions. At test time they are evaluated exactly on each image, and the number of satisfied
constraints decides how much of the proposed change is applied. That second mechanism is what makes
the constraints binding rather than advisory, and it is described in Section~\ref{sec:safety}.
